\chapter*{Abstract}
\initial{N}ational incident response teams face a structural imbalance: alert
volume grows faster than the workforce able to triage it, and response time
remains bounded by analyst availability. Existing orchestration platforms only
partly address this, because they stop at recommendation: the decision, and
therefore the waiting, stays human. This work asks the opposite question ---
under what conditions may a system act on its own? --- and answers it with a
single constraint: \emph{reversibility is the condition of autonomy}. The
resulting platform, CIRTDEFENSE, executes without prior validation only those
actions it knows how to undo, refuses to act on an undocumented context, applies
a natural-language policy compiled ahead of time, watches the effect of its own
actions and rolls them back by itself when they degrade the target. Faced with a
threat absent from its catalogue, it does not stay silent: it derives containment
from observed indicators alone, engages what is reversible and refers to a human
whatever leaves a durable trace. Every action is written to a hash-chained ledger,
the only opposable record. Modelling follows the unified process and the UML~2.5
notation; the implementation is validated by XXX automated tests, including XX
acceptance criteria derived directly from the requirements.

\paragraph{Keywords: autonomous incident response, reversibility, security
orchestration, auditability, UML.}

\clearpage
